\addtocontents{toc}{\protect\Needspace{3\baselineskip}}
\originalchapter{8.821 课程详讲：正则振子、零点能与引力子}
本章对应第 11--12 讲的谱与低能场. 目标是从横向正则对易关系读出开闭弦的态, 理解正规序和临界一致性, 再指出后续 D 膜构造需要的超弦场. 真正相互作用振幅、完整 BRST 同调和一般背景的 Weyl beta 函数另列高阶待补, 不把本章的自由谱计算当作那些工作已经完成.

\originalsection{开闭弦的零模系数为何不同}
共形规范的横向正则动量密度是 $\Pi^I=\dot X^I/(2\pi\alpha')$. 总动量 $p^I=\int\dd\sigma\,\Pi^I$ 应只由零模提供. 开弦区间长度为 $\pi$, 闭弦为 $2\pi$, 因而它们的线性零模分别是 $2\alpha'p^I\tau$ 和 $\alpha'p^I\tau$. 一套一致模式写作
\begin{align*}
X^I_{\rm open}&=x^I+2\alpha'p^I\tau
+\ii\sqrt{2\alpha'}\sum_{n\ne0}\frac{\alpha_n^I}{n}
\ee^{-\ii n\tau}\cos n\sigma,\\
X^I_{\rm closed}&=x^I+\alpha'p^I\tau
+\ii\sqrt{\frac{\alpha'}2}\sum_{n\ne0}\frac1n
\left[\alpha_n^I\ee^{-\ii n(\tau-\sigma)}+
\widetilde\alpha_n^I\ee^{-\ii n(\tau+\sigma)}\right].
\end{align*}
振子满足 $\alpha_n^\dagger=\alpha_{-n}$, 保证 $X^I$ Hermitian. 因为对时间微分抵消了 $\ii/n$, 这个因子同时控制坐标与动量的相对相位, 不能从模式式中任意省去.

令 $[X^I(\sigma),\Pi^J(\sigma')]=\ii\delta^{IJ}\delta(\sigma-\sigma')$, 用对应区间的 Fourier 完备关系得到
\[
[x^I,p^J]=\ii\delta^{IJ},\qquad
[\alpha_m^I,\alpha_n^J]=m\delta^{IJ}\delta_{m+n,0}.
\]
闭弦左、右两份有相同代数, 彼此对易. 对 $n>0$, $a_n^I=\alpha_n^I/\sqrt n$ 是湮灭算符, $a_n^{I\dagger}=\alpha_{-n}^I/\sqrt n$ 是创造算符. 波函数中的负 Fourier 指标与正能量创建并不矛盾, 能量由相应的数算符决定.

定义各模式占据数 $\mathcal N_n^I=\alpha_{-n}^I\alpha_n^I/n$ 和总层级
\[
N=\sum_{I,n>0}n\mathcal N_n^I
=\sum_{I,n>0}\alpha_{-n}^I\alpha_n^I,
\qquad \widetilde N\text{ 对左模同样定义}.
\]
一份 $n=2$ 振子把层级增加 2, 不等于简单数“创造算符有几个”. 例如 $\alpha_{-2}^I|0\rangle$ 和 $\alpha_{-1}^I\alpha_{-1}^J|0\rangle$ 都位于 $N=2$.

\originalsection{约束的零模给出质量与层级匹配}
动量满足 $M^2=-p^2=2p^+p^--p_I^2$. 将前章两个约束代入模式并取零 Fourier 分量, 开弦得到一套横向振子能量; 闭弦得到左右两套. 对闭弦积分 $X'^-$ 一圈必须为零, 于是左右总层级相等:
\[
N=\widetilde N.
\]
它是 $\sigma$ 周期一致性产生的约束, 不是左右每个 $n,I$ 占据数都必须分别相等. 左侧可用两个 $n=1$ 振子, 右侧用一个 $n=2$ 振子, 仍满足层级匹配.

正规序后, 用每一套横向玻色振子的截距 $a$ 记零点能, 质量公式为
\[
\alpha'M^2_{\rm open}=N-a,
\qquad
\alpha'M^2_{\rm closed}=2(N+\widetilde N-2a)=4(N-a).
\]
闭弦公式中最后一个等号使用层级匹配. 这套记法把截距取为正 $a$, 有的讲义把质量常数本身记为负数 $a_o,a_c$, 比较前需核对定义和 $\alpha'$ 的量纲.

\originalsection{发散级数的正规化不能当作普通求和}
每个横向自由玻色场的零点能是 $\tfrac12\sum_{n\ge1}n$. 这个级数在普通意义下发散, 并不等于 $-1/24$. 可用指数调节
\[
\sum_{n=1}^\infty n\ee^{-\epsilon n}
=\frac{\ee^{-\epsilon}}{(1-\ee^{-\epsilon})^2}
=\frac1{\epsilon^2}-\frac1{12}+O(\epsilon^2).
\]
移除与局部真空项有关的发散部分, 在与世界面一致性相容的正规化中有限部分是 $-1/12$. Zeta 解析延拓 $\zeta(-1)=-1/12$ 给同一有限项, 但其含义是正规化后的值, 不是声称普通正数相加结果为负. 横向 $D-2$ 个玻色场给
\[
E_0=-\frac{D-2}{24},\qquad a=\frac{D-2}{24}.
\]
因此开弦真空是 $M^2=-(D-2)/(24\alpha')$, 闭弦真空是 $M^2=-(D-2)/(6\alpha')$. 对 $D>2$ 是快子, 表示所选玻色真空有不稳定方向, 并非一个可在稳定 Minkowski 真空中任意传递超光速消息的普通粒子.

这一步不从正规化本身推出非微扰弦论已经定义完成. 一个形式渐近级数的解析延拓也不能唯一确定所有非微扰项. 谱的正规序常数需和 Lorentz/Weyl 一致性共同检查; 这是它在物理上被固定的条件.

\originalsection{临界维数与低层级表示检查}
第一开弦激发为 $\alpha_{-1}^I|0,p\rangle$, 有 $D-2$ 个横向分量, 且 $\alpha'M^2=1-(D-2)/24=(26-D)/24$. 无质量矢量在 $D$ 维有 $D-2$ 个物理偏振, 有质量矢量则有 $D-1$ 个. 因此把首激发解释为所需 Lorentz 矢量的必要一致性线索是让它无质量, 给 $D=26,a=1$.

这个表示计数是课程提供的直观检查, 完整的光锥 Lorentz boost 代数或共形 Weyl 异常消除才检验所有态的一致性. 本章使用标准平直自由玻色弦的临界结果, 不把一个层级的维数计数冒充整个量子 Lorentz 代数的逐项证明. PS3.2 给出横向 Virasoro 中心项的明确计算, 更完整规范固定鬼场/BRST 结构的范围另列.

在 $D=26$ 的 $N=2$ 开弦层级, $\alpha_{-1}^I\alpha_{-1}^J$ 对 $I,J$ 对称, 有 $24\times25/2=300$ 个态, $\alpha_{-2}^I$ 再给 24, 共 324. 有质量粒子 little group 为 $SO(25)$, 对称无迹二阶张量维数 $25\times26/2-1=324$, 恰好重组为完整有质量自旋-two 表示. 此例显示光锥 $SO(24)$ 下分开的两类振子怎样组成有质量 Lorentz 表示.

闭弦第一匹配层级 $N=\widetilde N=1$ 有 $\alpha_{-1}^I\widetilde\alpha_{-1}^J|0,p\rangle$, 共 $24^2=576$ 个无质量态. 分解为
\[
576=299+276+1,
\]
分别是对称无迹张量、反对称张量和迹. 它们对应目标场 $G_{\mu\nu}$ 的引力子偏振、二形式 $B_{\mu\nu}$、标量膨胀子 $\Phi$. 引力子的物理偏振数也由 $D(D-3)/2=26\times23/2=299$ 检查.

闭弦 $N=2$ 左右各有324态, 总数 $324^2=104976$. 它包含最高自旋四等多个有质量表示, 不是全层只有一个不可约“自旋-four”表示. 开闭弦质量平方层间隔分别为 $1/\alpha'$、$4/\alpha'$. 更高层的退化数由生成函数计数, 本卷 HW11.2--5 已给完整要求和算例.

\originalsection{低能场与正确的能量尺度}
弦重激发质量量级为 $M_s=1/\sqrt{\alpha'}$. 因此低能条件是 $E\sqrt{\alpha'}\ll1$, 不是把能量同 $1/\alpha'$ 这个质量平方直接比较. 在合适稳定背景和弱耦合下, 可积分掉重模式, 用无质量场的导数展开描述传播和散射. 玻色快子真空不稳定, 把它简单删除后的无质量作用量只是受限形式描写, 不能称为完整稳定玻色理论.

在一套常用弦参考系曲率约定下, NS--NS 无质量场领先作用量为
\[
S_{\rm NSNS}=\frac1{2\kappa_0^2}
\int\dd^DX\sqrt{-G_s}\,\ee^{-2\Phi}
\left[R_s+4G_s^{\mu\nu}\partial_\mu\Phi\partial_\nu\Phi
-\frac1{12}H_{\mu\nu\rho}H^{\mu\nu\rho}\right],
\]
其中 $H=\dd B$, 即 $H_{\mu\nu\rho}=\partial_\mu B_{\nu\rho}+\partial_\nu B_{\rho\mu}+\partial_\rho B_{\mu\nu}$. 这里把领先作用量作为已知弦振幅/世界面一致性结果引入; 课程未逐项算完所有 beta 函数, 本章也不声称仅靠偏振计数就已经推导全部相互作用系数.

为什么要区分 string 与爱因斯坦参考系度规? String frame 的 $R_s$ 前有场依赖因子 $\ee^{-2\Phi}$. 取
\[
G_{E\mu\nu}=\ee^{-4\Phi/(D-2)}G_{s\mu\nu}
\]
后, 利用 $\sqrt G$ 的 Weyl 权和曲率变换, $R_E$ 前的该因子抵消. 扔去合适边界条件下的全导数项, 爱因斯坦参考系形式的膨胀子动能是 $-4(\partial\Phi)^2/(D-2)$, 正常负时间号下具有正确能量. 对 $D=26$ 指数是 $-\Phi/6$, 对 $D=10$ 是 $-\Phi/2$; 直接代入指数可核对两种常见约定, 不能随意反号.

常数 $\Phi_0$ 令 $g_s=\ee^{\Phi_0}$, 改变引力有效耦合. $[\kappa_0^2]=L^{D-2}$, 因而任何具体 $\alpha'$ 幂次也必须随维数调整. 单凭 $\alpha'$ 的长度平方定义, 就能排除把十维、二十六维的 Newton 常数写成同一个无量纲公式.

\originalsection{超弦自由谱提供哪些替代场}
Ramond--Neveu--Schwarz（RNS）描述在世界面玻色场 $X^\mu$ 之外加入费米场 $\psi^\mu$, 适当规范固定后有八个横向玻色场和费米场. 费米场的空间边界可为 NS 半整数模式或 R 整数模式. 调节零点能时, 每个整数玻色场贡献 $-1/24$, 整数 R 费米场贡献 $+1/24$, 相消给 $a_R=0$. 半整数 NS 费米场的有限贡献为 $-1/48$, 与玻色场合为 $-1/16$, 八份给 $a_{\rm NS}=1/2$.

R 费米场的零模满足 Clifford 代数, 把真空变成 $SO(8)$ 旋量; Gliozzi--Scherk--Olive（GSO）投影选择适当手征性. NS 中选择保留相应奇费米场激发而移去真空快子, 最低态 $\psi_{-1/2}^I|0\rangle$ 给8维矢量. 因而单侧无质量 NS 和 R 各有8态. 使用10维临界超弦的一致性结果; 世界面超鬼场及完整 BRST 推导不是本课程这段自由谱介绍的已证明部分.

闭弦由左右两侧相乘. NS--NS 的 $8_v\otimes8_v=35+28+1$, 仍给 graviton、$B_2$、膨胀子, 共64个玻色偏振. NS--R 与 R--NS 共给128个费米偏振. R--R 因左右手征性选择分为:
\[
\begin{array}{c|c|c}
\text{理论}&\text{RR 势}&\text{物理 无质量 偏振数}\\\hline
\text{IIA（相反 手征性）}&C_1,C_3&8+56=64\\
\text{IIB（相同 手征性）}&C_0,C_2,C_4&1+28+35=64
\end{array}
\]
IIB 的四形式场强 $F_5$ 有 self-duality 条件, 才把四形式的70个横向分量减半为35. 这些场稍后让 D 膜能携带 RR 荷. 双侧总谱有128个玻色、128个费米无质量自由度, 与前面的单侧8态不同, 计数时应说明所在扇区.

IIA/IIB 的领先有效描述分别是对应 supergravity. 更广的超弦对偶网与十一维 M-theory 是高阶研究线索, 不由本章的自由谱计数建立其完整非微扰定义或唯一性. 一致背景中弦振幅的紫外性质也不能只由图的拓扑一句话证明: 世界面模空间退化、红外发散、快子和蝌蚪项仍要处理. 本课程接下来转向 D 膜这一实际构造工具, 保留这些未完成的理论边界.
